% Run pdfLaTeX from the source directory. Keep the lookup for the activity's year.
% Exactly ONE worked example, ONE guided example, and FIVE independent exercises.
% Counts apply to the whole activity; never repeat these blocks per family.
% Printable student activity only. Required gray headings and white criterion boxes
% are structural styles; semantic instructional highlights and legends are forbidden.
% Strip source highlighting while preserving content and ordinary formatting.
% Avoid inherited highlight-producing macros (e.g. the shaded preamble boxed macro).
\makeatletter
\def\input@path{{../../../preamble/}}
\makeatother
\input{preamble.tex}
\usepackage{needspace}
\usepackage{varwidth}

\geometry{letterpaper,margin=0.25in}
\pagestyle{empty}
\linespread{1.04}
\setlength{\parskip}{4pt}
\raggedbottom

% Exactly TWO table-style families. All helpers size to local \linewidth.
% Family 1: plain opening notation/formula reference, never a highlight legend.
% Adapt widths or use full width for the topic's actual definitions and formulas.
\newcommand{\NotationReferenceTable}[1]{%
    \par\begingroup
    \small
    \setlength{\tabcolsep}{3pt}
    \setlength{\extrarowheight}{0pt}
    \setlength{\parskip}{0pt}
    \renewcommand{\arraystretch}{1}
    \noindent\begin{tabularx}{\linewidth}{@{}>{\raggedright\arraybackslash}p{.21\linewidth}>{\raggedright\arraybackslash}p{.34\linewidth}>{\raggedright\arraybackslash}X@{}}
        \toprule
        \textbf{Symbol} & \textbf{Concept} & \textbf{Formula} \\
        \midrule
        #1
        \bottomrule
    \end{tabularx}\par
    \endgroup
}
% Family 2: all instructional/problem tables, in the printable activity.
% Width variants below are the SAME ruled family; adapt column count if needed.
% Body rows end with \\ \hline. Keep group/case headings short or wrap them.
\newcommand{\ClassroomRuledTable}[4]{%
    \par\begingroup
    \small
    \setlength{\tabcolsep}{4pt}
    \setlength{\extrarowheight}{0pt}
    \setlength{\parskip}{0pt}
    \renewcommand{\arraystretch}{1.16}
    \noindent\begin{tabularx}{\linewidth}{|>{\raggedright\arraybackslash}p{#1}|>{\raggedright\arraybackslash}p{#2}|>{\raggedright\arraybackslash}X|}
        \hline
        #3 \\ \hline
        #4
    \end{tabularx}\par
    \endgroup
}
\newcommand{\InformationTable}[1]{%
    \ClassroomRuledTable{.14\linewidth}{.32\linewidth}
        {\textbf{Symbol} & \textbf{Concept} & \textbf{In context...}}{#1}%
}
\NewDocumentCommand{\CalculationTable}{O{.10\linewidth} O{.20\linewidth} +m}{%
    \ClassroomRuledTable{#1}{#2}
        {\textbf{Symbol} & \textbf{Concept} & \textbf{Calculation, result, and meaning in context}}{#3}%
}
\newcommand{\TableGroup}[1]{\multicolumn{3}{|l|}{\textbf{#1}}\\ \hline}
% Structural styles for the printable activity only; no semantic highlighting.
\newcommand{\stageheading}[1]{%
    \par\Needspace{4\baselineskip}\addvspace{4pt}%
    {\centering\large\bfseries
     \tikz[baseline=(title.base)]{%
         \node[fill=black!10,text=black,rounded corners=2mm,
             inner sep=3pt,outer sep=0pt] (title) {\strut #1};}\par}%
    \nopagebreak[4]\smallskip
}
\newcommand{\criterionbox}[2]{%
    \begin{tcolorbox}[colback=white,colframe=black!45,boxrule=.4pt,
        arc=2mm,boxsep=0pt,left=6pt,right=6pt,top=5pt,bottom=5pt,
        before skip=4pt,after skip=4pt,width=\linewidth]
        \textbf{#1} #2
    \end{tcolorbox}%
}
\newcommand{\problemtitle}[3]{%
    \par\Needspace{4\baselineskip}\addvspace{0.4\baselineskip}%
    \noindent
    {\begin{varwidth}[t]{0.75\linewidth}%
         \bfseries\strut #1. #2%
     \end{varwidth}}%
    \hfill\textbf{#3}\par\nopagebreak[4]
}
\newcommand{\task}[2]{\par\noindent\textbf{#1} #2\par}
\newcommand{\responseline}{\par\noindent\rule{0pt}{1.55em}\hrulefill\par}
\newcommand{\workspace}[1]{\par\noindent\rule{0pt}{#1}\par\nopagebreak[4]}


% Local graphs use only the packages in the year-local preamble.
\newcommand{\EVGraph}[4]{%
\begin{tikzpicture}\begin{axis}[
 width=.88\linewidth,height=5.4cm,xmin=0,xmax=1,
 ymin={min(#2,#3)-.2*abs(#3-#2)},ymax={max(#2,#3)+.2*abs(#3-#2)},
 axis lines=left,xlabel={$p=P(E)$},ylabel={Expected net profit (USD)},
 title={$\mathbb E[Y_{#1}]$},xtick={0,1},ytick={#2,0,#3},
 extra x ticks={#4},extra x tick labels={$#4$},
 extra x tick style={tick label style={font=\small}},grid=major]
 \addplot[black,thick,domain=0:1,samples=2]{#2+(#3-#2)*x};
 \addplot[black,dashed,domain=0:1,samples=2]{0};
 \addplot[black,only marks,mark=*] coordinates {(#4,0)};
\end{axis}\end{tikzpicture}}
\newcommand{\BlankGraph}[1]{%
\begin{tikzpicture}[x=6.1cm,y=1cm]
 \draw[step=.61cm,black!20,very thin] (0,-1.4) grid (1,1.4);
 \draw[->] (0,-1.4)--(0,1.7) node[above,font=\small] {$\mathbb E[Y_{#1}]$ (USD)};
 \draw[->] (0,0)--(1.08,0) node[right,font=\small] {$p$};
 \node[below left] at (0,0) {$0$}; \node[below] at (1,0) {$1$};
\end{tikzpicture}}

% Short corresponding blocks can break naturally between paired procedures.
\NewDocumentCommand{\ColumnPair}{+m +m}{%
\par\noindent\begin{minipage}[t]{.49\linewidth}\vspace{0pt}\raggedright
#1
\end{minipage}\hfill
\begin{minipage}[t]{.49\linewidth}\vspace{0pt}\raggedright
#2
\end{minipage}\par}
\newcommand{\calcblank}{\underline{\hspace{1.15cm}}}
\newcommand{\shortblank}{\underline{\hspace{.7cm}}}
\newcommand{\PayoffSummary}[1]{%
\par\begingroup\small\setlength{\tabcolsep}{4pt}
\setlength{\parskip}{0pt}\renewcommand{\arraystretch}{1.16}
\begin{tabularx}{\linewidth}{|>{\raggedright\arraybackslash}p{.44\linewidth}|>{\raggedright\arraybackslash}X|}\hline
\textbf{Payoff summary} & \textbf{Value}\\\hline
#1\end{tabularx}\par\endgroup}
\begin{document}

{\centering {\Large\bfseries Bernoulli Models and Individual Break-Even Analysis\par}
Global Economics \textbar{} Grade 10 \textbar{} Term 2 \textbar{} Lesson 1\par}
\criterionbox{C3:}{Models binary economic outcomes with Bernoulli random variables, interpreting success and failure in contexts such as default, gain, loss, or claim occurrence.}

\stageheading{Equations and basic procedure}
One observation in its stated period has two mutually exclusive, exhaustive states.
The event label ``success'' means that $E$ occurs; it need not mean a monetary gain.
All data are hypothetical. Each $Y_i$ has one strictly positive and one strictly negative payoff.
\NotationReferenceTable{
$E, E^c$ & Event and complement & $E^c$: all outcomes outside $E$.\\
$X,p$ & Unitless Bernoulli indicator & $X=1$ on $E$, $X=0$ on $E^c$; $p=P(E)\in[0,1]$.\\
$1-p$ & Complementary mass & $P(X=0)=1-p$; $P(X=1)=p$.\\
$S_X,S_{Y_i}$ & Positive-probability supports & $S_X=\{x\in\mathbb R:P(X=x)>0\}$; $S_{Y_i}=\{y\in\mathbb R:P(Y_i=y)>0\}$.\\
$f_i,Y_i$ & State-to-money rule (USD) & $f_i:\{0,1\}\to\mathbb R$; $Y_i=f_i(X)$.\\
$f_i(0),f_i(1)$ & Fixed state payoffs (USD) & $Y_i=f_i(0)+[f_i(1)-f_i(0)]X$; $f_i(0)f_i(1)<0$.\\
$m_Z(z)$ & Probability mass function & $m_Z(z)=P(Z=z)$; zero outside the stated values.\\
$F_Z(t)$ & Cumulative distribution function & $F_Z(t)=P(Z\leq t)$; includes mass at the threshold.\\
$\mathbb E[Y_i]$ & Expected net profit (USD) & $p f_i(1)+(1-p)f_i(0)=f_i(0)+[f_i(1)-f_i(0)]p$.\\
$p_i^*$ & Individual break-even probability & $p_i^*=-f_i(0)/[f_i(1)-f_i(0)]\in(0,1)$.\\
}
\textbf{Indicator distribution (valid also at $p=0,1$):}
\[
m_X(x)=\begin{cases}1-p,&x=0,\\p,&x=1,\\0,&x\notin\{0,1\},\end{cases}
\qquad F_X(t)=\begin{cases}0,&t<0,\\1-p,&0\leq t<1,\\1,&t\geq1.\end{cases}
\]
\textbf{Monetary distribution:} write $a=\min(f_i(0),f_i(1))$, $b=\max(f_i(0),f_i(1))$ and $q=P(Y_i=a)$.
\[
m_{Y_i}(y)=\begin{cases}1-p,&y=f_i(0),\\p,&y=f_i(1),\\0,&\text{otherwise},\end{cases}
\qquad F_{Y_i}(t)=\begin{cases}0,&t<a,\\q,&a\leq t<b,\\1,&t\geq b.\end{cases}
\]
If $f_i(0)<0$, then $q=1-p$ and $\{Y_i<0\}=\{X=0\}$.
If $f_i(1)<0$, then $q=p$ and $\{Y_i<0\}=\{X=1\}$.
For $0<p<1$, $S_X=\{0,1\}$ and $S_{Y_i}=\{f_i(0),f_i(1)\}$.
At $p=0$ only $0$ and $f_i(0)$ remain in their respective supports; at $p=1$ only $1$ and $f_i(1)$ remain.
Zero-mass values cause no CDF jump.

\stageheading{General tasks}
\noindent\begin{minipage}[t]{.49\linewidth}\vspace{0pt}
\task{Task 1 --- Binary model and distributions.}{}
\begin{enumerate}[label=\arabic*.,itemsep=2pt,topsep=2pt,leftmargin=1.5em]
\item Observation, period, states $E,E^c$, shared unitless $X$, and monetary rules $Y_i=f_i(X)$ with state payoffs (USD).
\item Positive-probability set-builder supports, complete PMFs, and numerically ordered piecewise CDFs for $X$ and each $Y_i$, with numerical initial masses and supports at $p=0,1$.
\item Each loss event $\{Y_i<0\}$ in $X$, its probability function, initial probability, and contextual meaning.
\end{enumerate}
\end{minipage}\hfill
\begin{minipage}[t]{.49\linewidth}\vspace{0pt}
\task{Task 2 --- Individual expected values.}{}
\begin{enumerate}[label=\arabic*.,itemsep=2pt,topsep=2pt,leftmargin=1.5em]
\item Each expectation formula for fixed payoffs and varying $p\in[0,1]$, exact break-even probability, and separate positive and negative expectation intervals.
\item Extended payoff tables: state payoffs, expectation formulas, exact roots, and initial expectations (USD); contextual expected gain, loss, and break-even per period.
\item Separately labelled expected-value graphs: $p\in[0,1]$, monetary scale in USD, endpoint payoffs, expectation line, zero line, and exact break-even probability.
\end{enumerate}
\end{minipage}\par


\stageheading{Worked example}
\problemtitle{1}{Festival stall and quiet-day studio}{Complete}
A craft business observes one weekend at a time. Visitor demand is either high or low;
initially, high demand occurs on three out of every eight weekends. After all costs,
a festival stall earns a net profit of USD 105 with high demand and makes a loss of
USD 45 with low demand. A quiet-day studio programme makes a loss of USD 30 with high
demand and earns a net profit of USD 90 with low demand. Both activities face the same
weekend demand.

\task{Task 1 --- Establish the binary model and monetary transformations.}{}
\textbf{Shared model.} One observation is the business's outcome over one weekend.
Let $E$ be high demand and $E^c$ low demand; exactly one occurs.
Set $X=1$ on $E$, $X=0$ on $E^c$. Then $X\sim\operatorname{Bernoulli}(p)$ is unitless,
$p=P(E)$, initially $p=3/8$, and both activities use this same $X$; no independence is assumed.
For $0<p<1$,
\[
S_X=\{x\in\mathbb R:P(X=x)>0\}=\{0,1\},\quad
m_X(x)=\begin{cases}1-p,&x=0,\\p,&x=1,\\0,&\text{otherwise},\end{cases}\quad
F_X(t)=\begin{cases}0,&t<0,\\1-p,&0\leq t<1,\\1,&t\geq1.\end{cases}
\]
Initially $P(X=0)=5/8$, $P(X=1)=3/8$ and the CDF plateau is $5/8$.
At $p=0$, $S_X=\{0\}$; at $p=1$, $S_X=\{1\}$.

\ColumnPair{
\textbf{$Y_1$: festival stall net profit (USD).}\par
\textbf{State payoffs and affine rule.}
\[
\begin{aligned}
f_1(0)&=-45<0,& f_1(1)&=105>0,\\
Y_1&=f_1(0)+[f_1(1)-f_1(0)]X\\
&=-45+[105-(-45)]X\\
&=-45+150X.
\end{aligned}
\]
\textbf{Support} (USD; $0<p<1$).
\[
\begin{aligned}
S_{Y_1}&=\{y\in\mathbb R:P(Y_1=y)>0\}\\
&=\{-45,105\}.
\end{aligned}
\]
}{
\textbf{$Y_2$: studio net profit (USD).}\par
\textbf{State payoffs and affine rule.}
\[
\begin{aligned}
f_2(0)&=90>0,& f_2(1)&=-30<0,\\
Y_2&=f_2(0)+[f_2(1)-f_2(0)]X\\
&=90+[(-30)-90]X\\
&=90-120X.
\end{aligned}
\]
\textbf{Support} (USD; $0<p<1$).
\[
\begin{aligned}
S_{Y_2}&=\{y\in\mathbb R:P(Y_2=y)>0\}\\
&=\{-30,90\}.
\end{aligned}
\]
}
\ColumnPair{
\textbf{$Y_1$ --- Complete PMF and ordered CDF.}
\[
m_{Y_1}(y)=\begin{cases}1-p,&y=-45,\\p,&y=105,\\0,&\text{otherwise},\end{cases}
\]
\[
F_{Y_1}(t)=\begin{cases}0,&t<-45,\\1-p,&-45\leq t<105,\\1,&t\geq105.\end{cases}
\]
}{
\textbf{$Y_2$ --- Complete PMF and ordered CDF.}
\[
m_{Y_2}(y)=\begin{cases}p,&y=-30,\\1-p,&y=90,\\0,&\text{otherwise},\end{cases}
\]
\[
F_{Y_2}(t)=\begin{cases}0,&t<-30,\\p,&-30\leq t<90,\\1,&t\geq90.\end{cases}
\]
}
\ColumnPair{
\textbf{$Y_1$ --- Initial probabilities and loss.}
\[
m_{Y_1}(-45)=\tfrac58,\quad m_{Y_1}(105)=\tfrac38.
\]
\[
\{Y_1<0\}=\{X=0\}=E^c,\quad P(Y_1<0)=1-p.
\]
Initially, the loss probability is $5/8$: in the long run, five out of eight weekends produce a stall loss.
\par\textbf{Endpoint supports} (USD).
\[
p=0:\ S_{Y_1}=\{-45\};\quad p=1:\ S_{Y_1}=\{105\}.
\]
}{
\textbf{$Y_2$ --- Initial probabilities and loss.}
\[
m_{Y_2}(-30)=\tfrac38,\quad m_{Y_2}(90)=\tfrac58.
\]
\[
\{Y_2<0\}=\{X=1\}=E,\quad P(Y_2<0)=p.
\]
Initially, the loss probability is $3/8$: in the long run, three out of eight weekends produce a studio loss.
\par\textbf{Endpoint supports} (USD).
\[
p=0:\ S_{Y_2}=\{90\};\quad p=1:\ S_{Y_2}=\{-30\}.
\]
}

\task{Task 2 --- Analyse each expected value separately.}{}
\textbf{Common assumptions.} Retain the shared indicator from Task 1, hold both state
payoffs of each variable fixed, and vary only $p=P(E)\in[0,1]$.
\ColumnPair{
\textbf{$Y_1$ --- Expectation derivation} (USD).
\[
\begin{aligned}
\mathbb E[Y_1]&=p f_1(1)+(1-p)f_1(0)\\
&=105p-45(1-p)\\
&=105p-45+45p=-45+150p.
\end{aligned}
\]
\textbf{Break-even equation.}
\[
\begin{aligned}
\mathbb E[Y_1]=0&\iff-45+150p=0\\
&\iff150p=45\\
&\iff p=\tfrac{45}{150}=\tfrac3{10}.
\end{aligned}
\]
}{
\textbf{$Y_2$ --- Expectation derivation} (USD).
\[
\begin{aligned}
\mathbb E[Y_2]&=p f_2(1)+(1-p)f_2(0)\\
&=-30p+90(1-p)\\
&=-30p+90-90p=90-120p.
\end{aligned}
\]
\textbf{Break-even equation.}
\[
\begin{aligned}
\mathbb E[Y_2]=0&\iff90-120p=0\\
&\iff-120p=-90\\
&\iff p=\tfrac{-90}{-120}=\tfrac34.
\end{aligned}
\]
}
\ColumnPair{
\textbf{$Y_1$ --- Positive and negative expectations.}
\[
\begin{aligned}
\mathbb E[Y_1]>0&\iff150p>45\\
&\iff p>\tfrac3{10}\\
&\iff p\in(\tfrac3{10},1],\\[3pt]
\mathbb E[Y_1]<0&\iff150p<45\\
&\iff p<\tfrac3{10}\\
&\iff p\in[0,\tfrac3{10}).
\end{aligned}
\]
The stall has expected net gain above $3/10$, expected net loss below $3/10$, and break-even at $3/10$.
}{
\textbf{$Y_2$ --- Positive and negative expectations.}
\[
\begin{aligned}
\mathbb E[Y_2]>0&\iff-120p>-90\\
&\iff p<\tfrac34\\
&\iff p\in[0,\tfrac34),\\[3pt]
\mathbb E[Y_2]<0&\iff-120p<-90\\
&\iff p>\tfrac34\\
&\iff p\in(\tfrac34,1].
\end{aligned}
\]
Division by $-120$ reverses the inequalities. The studio has expected net gain below $3/4$, expected net loss above $3/4$, and break-even at $3/4$.
}
\ColumnPair{
\textbf{$Y_1$ --- Initial value and interpretation.}
\[
\begin{aligned}
\mathbb E[Y_1]\big|_{p=3/8}&=105(\tfrac38)-45(\tfrac58)\\
&=\tfrac{315-225}{8}=\tfrac{45}{4}\\
&=11.25\ \text{USD}.
\end{aligned}
\]
The long-run expected gain is USD 11.25 per weekend.
\PayoffSummary{$f_1(0),f_1(1)$ (USD)&$-45,105$\\\hline
$\mathbb E[Y_1]$ (USD)&$-45+150p$\\\hline
Exact $p_1^*$&$3/10$\\\hline
Initial value (USD)&$45/4=11.25$\\\hline}
}{
\textbf{$Y_2$ --- Initial value and interpretation.}
\[
\begin{aligned}
\mathbb E[Y_2]\big|_{p=3/8}&=-30(\tfrac38)+90(\tfrac58)\\
&=\tfrac{-90+450}{8}=\tfrac{360}{8}\\
&=45\ \text{USD}.
\end{aligned}
\]
The long-run expected gain is USD 45 per weekend.
\PayoffSummary{$f_2(0),f_2(1)$ (USD)&$90,-30$\\\hline
$\mathbb E[Y_2]$ (USD)&$90-120p$\\\hline
Exact $p_2^*$&$3/4$\\\hline
Initial value (USD)&$45$\\\hline}
}
\ColumnPair{
\textbf{$Y_1$ --- Graph.}\par\centering
\EVGraph{1}{-45}{105}{3/10}\par Exact zero: $p_1^*=3/10$.
}{
\textbf{$Y_2$ --- Graph.}\par\centering
\EVGraph{2}{90}{-30}{3/4}\par Exact zero: $p_2^*=3/4$.
}
Solid lines show expectations; dashed horizontal lines show zero; dots mark the individual roots.
\textbf{Worked checks.} The endpoint expectations are the state payoffs:
$\mathbb E[Y_1](0)=-45$, $\mathbb E[Y_1](1)=105$;
$\mathbb E[Y_2](0)=90$, $\mathbb E[Y_2](1)=-30$ (USD).
At the roots, $-45+150(3/10)=0$ and $90-120(3/4)=0$.
Positive initial expectation does not guarantee realized profit: each weekend gives one of the two state payoffs.
At break-even the long-run average is zero although each realized payoff is nonzero.

\stageheading{Guided Example}
\problemtitle{2}{Bookstall stock for a weekend fair}{Scaffolded}
A bookstall observes one weekend at a time. Demand is either high or low; initially,
high demand occurs on two out of every seven weekends. After all costs, a compact-stock
policy earns a net profit of USD 72 with high demand and makes a loss of USD 24 with low
demand. An expanded-stock policy earns a net profit of USD 126 with high demand and makes
a loss of USD 54 with low demand. Both policies face the same weekend demand.

\task{Task 1 --- Establish the binary model and monetary transformations.}{}
\textbf{Shared model.} One observation is the bookstall's outcome over one weekend.
Let $E$ be high demand and $E^c$ low demand; exactly one occurs.
Set $X=1$ on $E$, $X=0$ on $E^c$. Then $X\sim\operatorname{Bernoulli}(p)$ is unitless,
$p=P(E)$, initially $p=2/7$, and both policies use the same $X$; no independence is assumed.
Complete the numerical entries throughout the scaffold.
\[
S_X=\{x\in\mathbb R:P(X=x)>0\}=\{0,1\}\quad(0<p<1),\quad
m_X(x)=\begin{cases}1-p,&x=0,\\p,&x=1,\\0,&\text{otherwise},\end{cases}\quad
F_X(t)=\begin{cases}0,&t<0,\\1-p,&0\leq t<1,\\1,&t\geq1.\end{cases}
\]
Initially $P(X=1)=2/7$, $P(X=0)=1-2/7=\calcblank$;
the CDF plateau is $\calcblank$.
At $p=0$, $S_X=\{\shortblank\}$; at $p=1$, $S_X=\{\shortblank\}$.

\ColumnPair{
\textbf{$Y_1$: compact-stock net profit (USD).}\par
\textbf{State payoffs and affine rule.}
\[
\begin{aligned}
f_1(0)&=-24<0,& f_1(1)&=72>0,\\
Y_1&=f_1(0)+[f_1(1)-f_1(0)]X\\
&=-24+[72-(-24)]X\\
&=-24+\calcblank X.
\end{aligned}
\]
\textbf{Support} (USD; $0<p<1$).
\[
\begin{aligned}
S_{Y_1}&=\{y\in\mathbb R:P(Y_1=y)>0\}\\
&=\{\calcblank,\calcblank\}.
\end{aligned}
\]
}{
\textbf{$Y_2$: expanded-stock net profit (USD).}\par
\textbf{State payoffs and affine rule.}
\[
\begin{aligned}
f_2(0)&=-54<0,& f_2(1)&=126>0,\\
Y_2&=f_2(0)+[f_2(1)-f_2(0)]X\\
&=-54+[126-(-54)]X\\
&=-54+\calcblank X.
\end{aligned}
\]
\textbf{Support} (USD; $0<p<1$).
\[
\begin{aligned}
S_{Y_2}&=\{y\in\mathbb R:P(Y_2=y)>0\}\\
&=\{\calcblank,\calcblank\}.
\end{aligned}
\]
}
\ColumnPair{
\textbf{$Y_1$ --- Complete PMF and ordered CDF.}
\[
m_{Y_1}(y)=\begin{cases}1-p,&y=\calcblank,\\[5pt]p,&y=\calcblank,\\[5pt]0,&\text{otherwise},\end{cases}
\]
\[
F_{Y_1}(t)=\begin{cases}0,&t<\calcblank,\\[5pt]1-p,&\shortblank\leq t<\shortblank,\\[5pt]1,&t\geq\calcblank.\end{cases}
\]
}{
\textbf{$Y_2$ --- Complete PMF and ordered CDF.}
\[
m_{Y_2}(y)=\begin{cases}1-p,&y=\calcblank,\\[5pt]p,&y=\calcblank,\\[5pt]0,&\text{otherwise},\end{cases}
\]
\[
F_{Y_2}(t)=\begin{cases}0,&t<\calcblank,\\[5pt]1-p,&\shortblank\leq t<\shortblank,\\[5pt]1,&t\geq\calcblank.\end{cases}
\]
}
\ColumnPair{
\textbf{$Y_1$ --- Initial probabilities and loss.}
\[
\begin{aligned}
m_{Y_1}(-24)&=1-\tfrac27=\calcblank,\\[5pt]
m_{Y_1}(72)&=\calcblank.
\end{aligned}
\]
$\{Y_1<0\}=\{X=\shortblank\}=E^c$.\par
$P(Y_1<0)=1-p=1-2/7=\calcblank$ initially.\par
In the long run, $\shortblank$ out of seven weekends produce a compact-stock loss.
\par\textbf{Endpoint supports} (USD).
\[
\begin{aligned}
p=0:&\ S_{Y_1}=\{\calcblank\},\\[5pt]
p=1:&\ S_{Y_1}=\{\calcblank\}.
\end{aligned}
\]
}{
\textbf{$Y_2$ --- Initial probabilities and loss.}
\[
\begin{aligned}
m_{Y_2}(-54)&=1-\tfrac27=\calcblank,\\[5pt]
m_{Y_2}(126)&=\calcblank.
\end{aligned}
\]
$\{Y_2<0\}=\{X=\shortblank\}=E^c$.\par
$P(Y_2<0)=1-p=1-2/7=\calcblank$ initially.\par
In the long run, $\shortblank$ out of seven weekends produce an expanded-stock loss.
\par\textbf{Endpoint supports} (USD).
\[
\begin{aligned}
p=0:&\ S_{Y_2}=\{\calcblank\},\\[5pt]
p=1:&\ S_{Y_2}=\{\calcblank\}.
\end{aligned}
\]
}
\task{Task 2 --- Analyse each expected value separately.}{}
\textbf{Common assumptions.} Retain the shared indicator from Task 1, hold both state
payoffs of each variable fixed, and vary only $p=P(E)\in[0,1]$.
\ColumnPair{
\textbf{$Y_1$ --- Expectation derivation} (USD).
\[
\begin{aligned}
\mathbb E[Y_1]&=p f_1(1)+(1-p)f_1(0)\\
&=72p-24(1-p)\\
&=72p-24+24p\\
&=-24+\calcblank p.
\end{aligned}
\]
\textbf{Break-even equation.}
\[
\begin{aligned}
\mathbb E[Y_1]=0&\iff-24+[72-(-24)]p=0\\
&\iff\calcblank p=24\\[5pt]
&\iff p=\frac{24}{\shortblank}=\calcblank.
\end{aligned}
\]
}{
\textbf{$Y_2$ --- Expectation derivation} (USD).
\[
\begin{aligned}
\mathbb E[Y_2]&=p f_2(1)+(1-p)f_2(0)\\
&=126p-54(1-p)\\
&=126p-54+54p\\
&=-54+\calcblank p.
\end{aligned}
\]
\textbf{Break-even equation.}
\[
\begin{aligned}
\mathbb E[Y_2]=0&\iff-54+[126-(-54)]p=0\\
&\iff\calcblank p=54\\[5pt]
&\iff p=\frac{54}{\shortblank}=\calcblank.
\end{aligned}
\]
}
\ColumnPair{
\textbf{$Y_1$ --- Positive and negative expectations.}
\[
\begin{aligned}
\mathbb E[Y_1]>0&\iff[72-(-24)]p>24\\
&\iff p>\calcblank\\[5pt]
&\iff p\in(\calcblank,1],\\[5pt]
\mathbb E[Y_1]<0&\iff[72-(-24)]p<24\\
&\iff p<\calcblank\\[5pt]
&\iff p\in[0,\calcblank).
\end{aligned}
\]
Compact stock has expected net gain above $\calcblank$, expected net loss below $\calcblank$,
and break-even at $\calcblank$.
}{
\textbf{$Y_2$ --- Positive and negative expectations.}
\[
\begin{aligned}
\mathbb E[Y_2]>0&\iff[126-(-54)]p>54\\
&\iff p>\calcblank\\[5pt]
&\iff p\in(\calcblank,1],\\[5pt]
\mathbb E[Y_2]<0&\iff[126-(-54)]p<54\\
&\iff p<\calcblank\\[5pt]
&\iff p\in[0,\calcblank).
\end{aligned}
\]
Expanded stock has expected net gain above $\calcblank$, expected net loss below $\calcblank$,
and break-even at $\calcblank$.
}
\ColumnPair{
\textbf{$Y_1$ --- Initial value and interpretation.}
\[
\begin{aligned}
\mathbb E[Y_1]\big|_{p=2/7}&=72(\tfrac27)-24(1-\tfrac27)\\
&=\frac{144-24(\shortblank)}7\\[5pt]
&=\calcblank\ \text{USD}.
\end{aligned}
\]
The long-run expected \underline{\hspace{1.4cm}} (gain/loss) is USD $\calcblank$ per weekend.
\PayoffSummary{$f_1(0),f_1(1)$ (USD)&$-24,72$\\\hline
$\mathbb E[Y_1]$ (USD)&$-24+\shortblank p$\\[5pt]\hline
Exact $p_1^*$&$\calcblank$\\[5pt]\hline
Initial value (USD)&$\calcblank$\\[5pt]\hline}
}{
\textbf{$Y_2$ --- Initial value and interpretation.}
\[
\begin{aligned}
\mathbb E[Y_2]\big|_{p=2/7}&=126(\tfrac27)-54(1-\tfrac27)\\
&=\frac{252-54(\shortblank)}7\\[5pt]
&=\calcblank\ \text{USD}.
\end{aligned}
\]
The long-run expected \underline{\hspace{1.4cm}} (gain/loss) is USD $\calcblank$ per weekend.
\PayoffSummary{$f_2(0),f_2(1)$ (USD)&$-54,126$\\\hline
$\mathbb E[Y_2]$ (USD)&$-54+\shortblank p$\\[5pt]\hline
Exact $p_2^*$&$\calcblank$\\[5pt]\hline
Initial value (USD)&$\calcblank$\\[5pt]\hline}
}
\ColumnPair{
\textbf{$Y_1$ --- Graph.}\par
Use a labelled USD scale; plot the endpoint payoffs, join them with the expectation line, and mark the exact root.
\par\centering\BlankGraph{1}\par Exact zero: $p_1^*=\calcblank$.
}{
\textbf{$Y_2$ --- Graph.}\par
Use a labelled USD scale; plot the endpoint payoffs, join them with the expectation line, and mark the exact root.
\par\centering\BlankGraph{2}\par Exact zero: $p_2^*=\calcblank$.
}
The horizontal axes provide the zero lines; $p=P(E)$ ranges from 0 to 1.
At either break-even probability, the long-run expected net profit is zero, although a realized weekend payoff is nonzero.

\stageheading{Independent Practice}
\problemtitle{3}{Bicycle-repair service}{Core: service}
A bicycle workshop observes one day at a time. Each day is either busy or quiet;
initially, busy days occur on three out of every five days. After all costs,
walk-in repairs earn a net profit of USD 84 on a busy day and make a loss of USD 36
on a quiet day. Quiet-day lessons make a loss of USD 48 on a busy day and earn a net
profit of USD 72 on a quiet day. Both policies face the same daily demand.

\problemtitle{4}{Micro-hydro electricity sales}{Core: sales}
A micro-hydro producer observes one month at a time. River flow is either strong or weak;
initially, strong flow occurs in two out of every five months. After all costs,
market sales earn a net profit of USD 98 with strong flow and make a loss of USD 42
with weak flow. A fixed-delivery arrangement earns a net profit of USD 70 with strong
flow and makes a loss of USD 70 with weak flow. Both arrangements face the same monthly flow.

\problemtitle{5}{Digital printing capacity}{Core: capacity}
A digital-printing business observes one quarter at a time. A bulk order either arrives
or does not arrive; initially, one arrives in four out of every nine quarters.
After all quarterly costs, the current printer earns a net profit of USD 150 with a bulk
order and makes a loss of USD 75 without one. A leased printer earns a net profit of
USD 375 with a bulk order and makes a loss of USD 75 without one. Both policies face
the same quarterly order outcome.

\problemtitle{6}{Courier disruption and backup}{Extension: precision}
A courier depot observes one week at a time. A route disruption either occurs or does not
occur; initially, one occurs in five out of every twelve weeks. After all costs, a standard
reserve policy makes a loss of USD 27.50 with disruption and earns a net profit of USD 82.50
without disruption. A flexible reserve policy makes a loss of USD 45 with disruption and
earns a net profit of USD 67.50 without disruption. Both policies face the same weekly
route outcome; the decimal amounts are exact.

\problemtitle{7}{Cold-storage service and event labels}{Extension: relabelling}
A cold-storage business observes one month at a time. Supply is either interrupted or
uninterrupted; initially, an interruption occurs in seven out of every twenty months.
After all costs, regular storage makes a loss of USD 64 with interruption and earns a net
profit of USD 96 with uninterrupted supply. Emergency storage earns a net profit of USD 112
with interruption and makes a loss of USD 28 with uninterrupted supply. Both arrangements
face the same monthly supply outcome. For each arrangement, also express net profit as a
constant plus a coefficient times an indicator that equals one for uninterrupted supply
and zero for interruption. Express its expected net profit in terms of the probability of
uninterrupted supply, which is one minus the interruption probability; find the exact
break-even probability of uninterrupted supply and calculate its initial expected net profit
under this relabelling. Use the interruption probability on the horizontal axis of both
expected-value graphs.

\end{document}
